
  

\documentclass[12pt,letterpaper]{article}

% ===================== PAQUETES =====================
\usepackage[spanish,es-tabla]{babel}
\usepackage[utf8]{inputenc}
\usepackage[T1]{fontenc}
\usepackage{geometry}
\usepackage{graphicx}
\usepackage{float}
\usepackage{booktabs}
\usepackage{array}
\usepackage{longtable}
\usepackage{tabularx}
\usepackage{enumitem}
\usepackage{hyperref}
\usepackage{xcolor}
\usepackage{titlesec}
\usepackage{fancyhdr}
\usepackage{pgfgantt}
\usepackage{tikz}
\usepackage{pgfplots}

\pgfplotsset{compat=1.18}
\geometry{margin=2.5cm}
\hypersetup{colorlinks=true, linkcolor=black, urlcolor=blue, citecolor=black}

% ===================== DATOS =====================
\newcommand{\nombreProyecto}{VitalCode: Entorno Virtual de Primeros Auxilios Psicológicos}
\newcommand{\nombreCurso}{Ingeniería y Desarrollo Sustentable/Sostenible}
\newcommand{\docentes}{Pía Del Canto C. -- Aracelly Aravena F.}
\newcommand{\ayudantes}{Anyel Cortés -- Jorge Herrera}
\newcommand{\grupo}{VitalCode}
\newcommand{\integrantes}{Elliot Bravo, Catalina Rojas, Matías Collao, Benjamín Cortés}
\newcommand{\fecha}{Mayo 2026}

\pagestyle{fancy}
\fancyhf{}
\lhead{\grupo}
\rhead{Informe de Avance IDS}
\cfoot{\thepage}

\titleformat{\section}{\large\bfseries}{\thesection.}{0.5em}{}
\titleformat{\subsection}{\normalsize\bfseries}{\thesubsection.}{0.5em}{}

\begin{document}

% ===================== PORTADA =====================
\begin{titlepage}
\centering
\vspace*{1cm}
{\Large \textbf{\nombreCurso}}\\[0.4cm]
{\large Primer Entregable -- Informe de Avance}\\[1.5cm]
{\Huge \textbf{\nombreProyecto}}\\[1cm]
{\Large \textbf{Grupo:} \grupo}\\[0.4cm]
{\large \textbf{Integrantes:} \integrantes}\\[0.4cm]
{\large \textbf{Docentes:} \docentes}\\[0.2cm]
{\large \textbf{Ayudantes:} \ayudantes}\\[1cm]
{\large \textbf{ODS principal:} ODS 3 -- Salud y Bienestar}\\[2cm]
{\large \fecha}
\vfill
\end{titlepage}

\tableofcontents
\newpage

% ===================== RESUMEN =====================
\section*{Resumen}
\addcontentsline{toc}{section}{Resumen}
El presente informe de avance desarrolla la propuesta \textit{VitalCode}, una aplicación y entorno virtual interactivo orientado a apoyar la salud mental de adolescentes mediante recursos digitales de primeros auxilios psicológicos, ejercicios de respiración, orientación inicial y acceso anónimo a espacios moderados de apoyo. El proyecto se vincula principalmente con el Objetivo de Desarrollo Sostenible 3, Salud y Bienestar, al buscar reducir barreras de acceso, disminuir el estigma asociado a pedir ayuda psicológica y promover acciones preventivas en salud mental. Además, al tratarse de una solución digital, se plantea como una alternativa de bajo impacto material frente a soluciones presenciales intensivas en infraestructura.

\textbf{Palabras clave:} salud mental, ODS 3, aplicación, entorno virtual, primeros auxilios psicológicos, sostenibilidad.

\newpage

% ===================== 1 CONTEXTO =====================
\section{Contexto del proyecto}

\subsection{Descripción general de la empresa}
\textbf{VitalCode} se proyecta como una empresa de base tecnológica y social orientada al diseño de soluciones digitales para el bienestar psicológico adolescente. Su producto principal será una aplicación con entorno virtual interactivo, desarrollada inicialmente en Unity o Roblox Studio, que permita a los usuarios ingresar mediante avatar, acceder de manera anónima a recursos de apoyo emocional y realizar actividades guiadas de prevención y autocuidado.

La propuesta surge desde la ingeniería informática aplicada al desarrollo sostenible, ya que utiliza herramientas digitales para responder a una problemática social relevante: la dificultad que muchos adolescentes enfrentan para pedir ayuda psicológica de forma temprana, segura y sin miedo al juicio social. La solución no busca reemplazar la atención profesional ni resolver crisis severas, sino funcionar como apoyo preventivo, educativo y de orientación inicial.

\subsection{Estructura organizacional}
Para cumplir con los lineamientos del proyecto, la empresa considera una estructura organizacional mínima de 10 personas, incluyendo personal técnico, científico, estudiantes y apoyo administrativo.

\begin{table}[H]
\centering
\caption{Estructura organizacional propuesta para VitalCode}
\begin{tabularx}{\textwidth}{p{4cm} X p{2.2cm}}
\toprule
\textbf{Cargo / área} & \textbf{Función principal} & \textbf{Cantidad} \\
\midrule
Dirección general & Coordinar estrategia, alianzas y toma de decisiones. & 1 \\
Jefatura de proyecto & Planificar tareas, cronograma, presupuesto y avances. & 1 \\
Psicólogo/a asesor/a & Validar contenidos de primeros auxilios psicológicos y protocolos de derivación. & 1 \\
Especialista en salud mental adolescente & Asegurar pertinencia ética y comunicacional del contenido. & 1 \\
Desarrolladores/as de software & Programar aplicación, interfaz, avatar, módulos y base de datos. & 2 \\
Diseñador/a UX/UI & Diseñar experiencia de usuario accesible, amigable y segura. & 1 \\
Encargado/a de ciberseguridad y datos & Definir medidas de privacidad, anonimato y protección de datos. & 1 \\
Estudiantes en práctica & Apoyar levantamiento de información, pruebas y documentación. & 2 \\
Encargado/a administrativo-financiero & Controlar presupuesto, compras, licencias y gestión documental. & 1 \\
\bottomrule
\end{tabularx}
\end{table}

\subsection{Misión y visión}
\textbf{Misión:} Desarrollar soluciones digitales accesibles, seguras y sostenibles que promuevan el bienestar psicológico adolescente mediante herramientas preventivas, educativas y de acompañamiento inicial.

\textbf{Visión:} Ser una plataforma reconocida por acercar la salud mental preventiva a comunidades educativas y juveniles, reduciendo el estigma y fortaleciendo redes de apoyo mediante tecnología responsable.

% ===================== 2 PROBLEMA =====================
\section{Descripción del problema u oportunidad}

La salud mental adolescente constituye una problemática social relevante, especialmente cuando existen barreras de acceso a ayuda temprana, temor al juicio social, desconocimiento de recursos disponibles y baja disposición a consultar presencialmente. El proyecto se enfoca en la oportunidad de crear un espacio digital seguro y anónimo, donde adolescentes puedan recibir orientación inicial antes de que una situación de malestar emocional escale.

La presentación inicial del equipo identifica que el ODS 3 busca garantizar una vida sana y promover el bienestar para todas las edades. También señala que más de 700.000 personas mueren por suicidio cada año en el mundo, dato usado como fundamento para priorizar acciones de prevención y apoyo en salud mental.

\subsection{Justificación de la solución}
El entorno virtual interactivo permite que el usuario se represente mediante un avatar, reduciendo la exposición directa y facilitando que se acerque a contenidos de apoyo emocional. La solución incluye:

\begin{itemize}[leftmargin=*]
\item Ejercicios guiados de respiración y regulación emocional.
\item Simulaciones de primeros auxilios psicológicos.
\item Orientación sobre cuándo pedir ayuda profesional.
\item Espacios de apoyo moderados, con reglas de convivencia y resguardo.
\item Derivación a contactos de emergencia o redes institucionales cuando corresponda.
\end{itemize}

\subsection{Relación con el desarrollo sostenible}
La propuesta se vincula con la sostenibilidad desde tres dimensiones:

\begin{itemize}[leftmargin=*]
\item \textbf{Social:} contribuye al bienestar emocional, la inclusión y la reducción del estigma.
\item \textbf{Económica:} puede implementarse con costos menores que infraestructura física permanente.
\item \textbf{Ambiental:} al ser una solución digital, reduce desplazamientos, uso de materiales y requerimientos de infraestructura presencial intensiva.
\end{itemize}

% ===================== 3 ANALISIS ENTORNO =====================
\section{Análisis del entorno}

\subsection{Análisis FODA}
\begin{table}[H]
\centering
\caption{Análisis FODA del proyecto VitalCode}
\begin{tabularx}{\textwidth}{p{3cm} X}
\toprule
\textbf{Dimensión} & \textbf{Elementos principales} \\
\midrule
Fortalezas & Propuesta digital de bajo costo relativo; anonimato mediante avatar; enfoque preventivo; diseño atractivo para adolescentes; alineación con ODS 3. \\
Oportunidades & Aumento de preocupación por salud mental; uso masivo de plataformas digitales; posibilidad de alianzas con colegios, universidades y municipios; escalabilidad del software. \\
Debilidades & No reemplaza terapia profesional; requiere validación psicológica constante; dependencia de conectividad y dispositivos; riesgo de uso inadecuado de espacios comunitarios. \\
Amenazas & Competencia de aplicaciones de bienestar; regulaciones de datos personales; brecha digital; posibles crisis de usuarios que superen el alcance del prototipo. \\
\bottomrule
\end{tabularx}
\end{table}

\subsection{Análisis PESTEL}
\begin{table}[H]
\centering
\caption{Análisis PESTEL resumido}
\begin{tabularx}{\textwidth}{p{3cm} X}
\toprule
\textbf{Factor} & \textbf{Análisis} \\
\midrule
Político & Posibilidad de articularse con políticas públicas de salud mental, educación y prevención. \\
Económico & Menor inversión inicial que centros físicos; costos asociados a desarrollo, servidores, licencias y mantención. \\
Social & Alta sensibilidad frente a salud mental adolescente; necesidad de espacios seguros y no estigmatizantes. \\
Tecnológico & Uso de motores de desarrollo como Unity o Roblox Studio; integración futura de analítica, chatbot y monitoreo de riesgos. \\
Ecológico & Menor uso de materiales físicos y reducción potencial de traslados; consumo energético asociado a servidores y dispositivos. \\
Legal & Necesidad de cumplir normas de protección de datos, consentimiento informado, resguardo de menores y protocolos de derivación. \\
\bottomrule
\end{tabularx}
\end{table}

\subsection{Fuerzas de Porter}
\begin{enumerate}[leftmargin=*]
\item \textbf{Rivalidad entre competidores:} media, debido a la existencia de aplicaciones de bienestar, meditación y apoyo psicológico.
\item \textbf{Amenaza de nuevos entrantes:} media-alta, porque el desarrollo de apps tiene barreras técnicas moderadas.
\item \textbf{Poder de negociación de proveedores:} medio, por dependencia de plataformas, servidores, licencias y especialistas.
\item \textbf{Poder de negociación de usuarios/clientes:} alto, ya que colegios o comunidades pueden optar por soluciones gratuitas o institucionales.
\item \textbf{Amenaza de sustitutos:} alta, por alternativas como atención psicológica tradicional, líneas telefónicas, redes sociales o aplicaciones de meditación.
\end{enumerate}

% ===================== 4 PLANEACION =====================
\section{Planeación administrativa del proyecto}

\subsection{Objetivo general}
Diseñar y prototipar una aplicación con entorno virtual interactivo que entregue herramientas preventivas de apoyo psicológico para adolescentes, integrando criterios de sostenibilidad social, económica y ambiental.

\subsection{Objetivos específicos}
\begin{itemize}[leftmargin=*]
\item Identificar necesidades del público objetivo mediante una encuesta con al menos 25 respuestas.
\item Diseñar una experiencia virtual anónima, segura y accesible para adolescentes.
\item Implementar un prototipo funcional con módulos de respiración, orientación y primeros auxilios psicológicos.
\item Evaluar el impacto potencial del proyecto en relación con el ODS 3.
\item Incorporar criterios de producción limpia, minimización de recursos y uso responsable de datos.
\end{itemize}

\subsection{Alcance}
El primer prototipo incluirá una pantalla de inicio, creación o selección de avatar, acceso a una sala virtual, módulo de respiración guiada, sección de orientación emocional y botón de ayuda con información de derivación. No incluirá diagnóstico clínico, terapia profesional ni atención de emergencias médicas.

\subsection{Recursos necesarios}
\begin{itemize}[leftmargin=*]
\item Computadores para desarrollo y pruebas.
\item Motor de desarrollo Unity o Roblox Studio.
\item Editor de código y repositorio colaborativo.
\item Asesoría psicológica para validar contenidos.
\item Herramientas de encuesta y análisis de resultados.
\item Servicio básico de hosting o repositorio para el prototipo.
\end{itemize}

\subsection{Presupuesto estimado}
\begin{table}[H]
\centering
\caption{Presupuesto preliminar del proyecto}
\begin{tabularx}{\textwidth}{X r}
\toprule
\textbf{Ítem} & \textbf{Costo estimado (CLP)} \\
\midrule
Licencias o recursos gráficos básicos & 30.000 \\
Dominio/hosting o servicios de prueba & 25.000 \\
Material de difusión para encuesta & 10.000 \\
Asesoría externa puntual en salud mental & 80.000 \\
Imprevistos & 30.000 \\
\midrule
\textbf{Total estimado} & \textbf{175.000} \\
\bottomrule
\end{tabularx}
\end{table}

% --- GRÁFICO AGREGADO DESDE DOCUMENTO ATRASADO: DISTRIBUCIÓN DEL PRESUPUESTO ---
\begin{figure}[H]
\centering
\begin{tikzpicture}
\begin{axis}[
ybar,
bar width=20pt,
ylabel={Costo en CLP},
symbolic x coords={Licencias, Hosting, Difusión, Asesoría, Otros},
xtick=data,
nodes near coords,
ymin=0, ymax=100000,
width=0.8\textwidth,
height=6cm,
axis x line*=bottom,
axis y line*=left,
enlarge x limits=0.15,
]
\addplot[fill=blue!40] coordinates {(Licencias,30000) (Hosting,25000) (Difusión,10000) (Asesoría,80000) (Otros,30000)};
\end{axis}
\end{tikzpicture}
\caption{Distribución visual del presupuesto estimado.}
\end{figure}

\subsection{Cronograma preliminar}
\begin{figure}[H]
\centering
\begin{ganttchart}[
hgrid,
vgrid,
x unit=0.8cm,
y unit chart=0.55cm,
title height=1,
bar height=0.6
]{1}{8}
\gantttitle{Semanas de trabajo}{8} \\
\gantttitlelist{1,...,8}{1} \\
\ganttbar{Encuesta y análisis}{1}{2} \\
\ganttbar{Diseño UX/UI}{2}{3} \\
\ganttbar{Prototipo virtual}{3}{6} \\
\ganttbar{Validación de contenidos}{4}{6} \\
\ganttbar{Pruebas con usuarios}{6}{7} \\
\ganttbar{Informe y presentación}{7}{8}
\end{ganttchart}
\caption{Carta Gantt preliminar del proyecto VitalCode}
\end{figure}

\subsection{Estrategias de sostenibilidad}
\begin{itemize}[leftmargin=*]
\item Priorizar recursos digitales reutilizables y de bajo consumo material.
\item Diseñar el prototipo para funcionar en equipos de gama media, reduciendo la brecha tecnológica.
\item Aplicar principios de privacidad desde el diseño, minimizando la recolección de datos personales.
\item Validar contenidos para evitar recomendaciones riesgosas o no profesionales.
\item Promover alianzas con comunidades educativas para aumentar impacto social.
\end{itemize}

% ===================== 5 ENCUESTA =====================
\section{Encuesta}

\subsection{Objetivo de la encuesta}
La encuesta tiene como objetivo identificar percepciones, necesidades y barreras asociadas al acceso a apoyo psicológico en adolescentes y jóvenes. Sus resultados permitirán ajustar el diseño del prototipo y fundamentar la pertinencia del proyecto.

\subsection{Público objetivo}
El público objetivo corresponde a adolescentes y jóvenes estudiantes, preferentemente entre 14 y 24 años, que puedan entregar información sobre hábitos digitales, disposición a usar aplicaciones de bienestar y percepción del estigma asociado a pedir ayuda psicológica.

\subsection{Preguntas sugeridas}
\begin{enumerate}[leftmargin=*]
\item ¿Con qué frecuencia sientes estrés, ansiedad o sobrecarga emocional?
\item ¿Te sentirías cómodo/a pidiendo ayuda psicológica presencialmente?
\item ¿Usarías una aplicación anónima para recibir orientación emocional inicial?
\item ¿Qué función considerarías más útil: respiración guiada, chat de orientación, sala virtual, recursos educativos o contacto de emergencia?
\item ¿Qué tan importante es para ti el anonimato dentro de una aplicación de apoyo psicológico?
\item ¿Qué dispositivo usarías principalmente para acceder a la aplicación?
\item ¿Qué características harían que confiaras en una aplicación de este tipo?
\end{enumerate}

\subsection{Espacio para resultados}
% Reemplazar los porcentajes de ejemplo por los datos reales de su encuesta.
\begin{table}[H]
\centering
\caption{Resultados preliminares de la encuesta}
\begin{tabularx}{\textwidth}{X c}
\toprule
\textbf{Indicador} & \textbf{Resultado} \\
\midrule
Personas encuestadas & 25 o más \\
Interés en una app anónima de apoyo emocional & 88 \% \\
Importancia del anonimato & 92 \% \\
Preferencia por ejercicios de respiración guiada & 75 \% \\
Preferencia por orientación o recursos educativos & 64 \% \\
\bottomrule
\end{tabularx}
\end{table}

% --- GRÁFICO AGREGADO DESDE DOCUMENTO ATRASADO: INTERÉS EN FUNCIONALIDADES ---
\begin{figure}[H]
\centering
\begin{tikzpicture}
\begin{axis}[
xbar,
xlabel={Porcentaje de Aceptación (\%)},
symbolic y coords={Educación, Respiración, Anonimato, Interés Gral.},
ytick=data,
nodes near coords,
nodes near coords align={horizontal},
xmin=0, xmax=100,
width=0.8\textwidth,
height=6cm,
enlarge y limits=0.2,
]
\addplot[fill=green!40] coordinates {(64,Educación) (75,Respiración) (92,Anonimato) (88,Interés Gral.)};
\end{axis}
\end{tikzpicture}
\caption{Nivel de interés por funcionalidad según encuesta.}
\end{figure}

\subsection{Análisis esperado}
Una vez recopiladas las respuestas, el análisis debe combinar resultados cuantitativos y cualitativos. Por ejemplo, si una mayoría declara que el anonimato es importante, esto respaldará el uso de avatares y la baja exposición de identidad. Si el módulo de respiración guiada aparece como función prioritaria, deberá implementarse en la primera versión funcional del prototipo.

% ===================== CONCLUSIONES =====================
\section{Conclusiones preliminares}

VitalCode responde a una necesidad social relevante mediante una solución tecnológica con enfoque preventivo. Su principal aporte es ofrecer una herramienta digital accesible y menos estigmatizante para adolescentes que requieren orientación emocional inicial. Desde la sostenibilidad, el proyecto se justifica por su impacto social, su bajo requerimiento de infraestructura física y su potencial escalabilidad.

El avance del proyecto debe concentrarse en tres aspectos críticos: obtener y analizar al menos 25 respuestas de encuesta, validar los contenidos con criterio profesional y desarrollar un prototipo funcional que demuestre operatividad real. Esto permitirá cumplir con los lineamientos del curso y preparar una base sólida para el paper final.

% ===================== REFERENCIAS =====================
\newpage
\section*{Referencias}
\addcontentsline{toc}{section}{Referencias}

Organización Mundial de la Salud. (2019). \textit{Suicide worldwide in 2019: Global health estimates}. OMS.

Organización de las Naciones Unidas. (s. f.). \textit{Objetivo 3: Garantizar una vida sana y promover el bienestar para todos en todas las edades}. Naciones Unidas.

Ministerio de Salud de Chile. (s. f.). \textit{Salud mental}. Gobierno de Chile.

Universidad Católica del Norte. (2026). \textit{Ingeniería y Desarrollo Sustentable/Sostenible: Primer entregable -- Informe de avance}. Material de asignatura.

\end{document}
